\documentclass[11pt]{amsart}
\usepackage[margin=1.04in]{geometry}
\usepackage[T1]{fontenc}
\usepackage{lmodern,amsmath,amssymb,amsfonts,mathtools,mathrsfs,booktabs,longtable,array,microtype}
\usepackage[hidelinks]{hyperref}
\ifdefined\pdfinfoomitdate\pdfinfoomitdate=1\fi
\ifdefined\pdftrailerid\pdftrailerid{}\fi
\ifdefined\pdfsuppressptexinfo\pdfsuppressptexinfo=15\fi
\newtheorem{theorem}{Theorem}[section]
\newtheorem{proposition}[theorem]{Proposition}
\newtheorem{lemma}[theorem]{Lemma}
\newtheorem{corollary}[theorem]{Corollary}
\theoremstyle{definition}\newtheorem{definition}[theorem]{Definition}
\newtheorem{example}[theorem]{Example}
\theoremstyle{remark}\newtheorem{remark}[theorem]{Remark}
\newcommand{\E}{\mathbb E}
\newcommand{\Pp}{\mathbb P}
\newcommand{\R}{\mathbb R}
\newcommand{\N}{\mathbb N}
\newcommand{\one}{\mathbf 1}
\newcommand{\osc}{\operatorname{osc}}
\newcommand{\TV}{\operatorname{TV}}
\newcommand{\Lip}{\operatorname{Lip}}
\newcommand{\conv}{\operatorname{conv}}
\newcommand{\diag}{\operatorname{diag}}
\newcommand{\rank}{\operatorname{rank}}
\newcommand{\tr}{\operatorname{tr}}
\newcommand{\supp}{\operatorname{supp}}
\newcommand{\dist}{\operatorname{dist}}
\newcommand{\calR}{\mathcal R}
\newcommand{\calC}{\mathcal C}
\newcommand{\calE}{\mathfrak E}
\newcommand{\calG}{\mathcal G}
\newcommand{\calK}{\mathcal K}
\newcommand{\norm}[1]{\left\lVert#1\right\rVert}
\newcommand{\abs}[1]{\left|#1\right|}
\DeclareMathOperator*{\esssup}{ess\,sup}
\title[General Theta Foundations I]{General Theta Foundations I:\\Acquired Geometry and Causal Resource Transfer}
\author{Qian Qi}
\date{10 October 2026}
\subjclass[2020]{Primary 60J20, 93E20; Secondary 62B15, 60K15, 94A17, 81P45}
\keywords{causal experiment, predictive quotient, acquired occupation, finite memory, Markov renewal, task-sensitive regularity, singular observation}
\begin{document}
\begin{abstract}
A finite observer may retain information across physical preparations even when the induced boundary chain approaches a change of recurrent rank. We characterize exactly when such degeneration is invisible to a specified bounded task. On a compact family with continuous finite responses and uniformly integrable physical return times, continuity of the physical-time gain is equivalent to uniform acquisition and discount limits and to uniformly bounded approximate marked Poisson correctors. The corrector uses the joint duration--exit matrix, not the product of mean duration and exit probability. It yields a same-memory quantitative transfer law and robust attainment without a cost-universal recurrence bound whenever the task criterion holds on the full declared program class. Applied to executable predictive tests, it controls the actual acquired occupation law; small-ball mass, global covers and occupation-weighted update stability then give matching checkpoint and causal memory bounds. For finite algebraic instruments, the criterion is decidable and admits algebraic convergence moduli; a separate result makes recurrence exhaustion quantitatively operational even when continuity fails. A rare-revelation experiment attains the corrector rate against all history-dependent strategies. A noncommuting quantum realization crosses a preparation-rank singularity, while controlled queue excursions give an unbounded physical-state realization. A common scored experiment has a matching acquisition--memory--calibration--deficiency law governed by acquired quantization and a task-visible modulus rather than ambient parameter dimension.
\end{abstract}
\maketitle
\section{Introduction}
A prepared experiment determines which future tests are executable, which histories have the same predictive content, and which distributions on that content are actually obtained. A finite observer must realize its predictions by one causal program. The geometry of a formal parameter space does not determine the cost of doing so. Neither does the rate at which an unrelated, cost-universal Markov chain mixes. The relevant obstruction is the portion of temporal degeneration that is visible to the task being performed.

This distinction becomes essential when physical preparation leaves the observer's memory intact. The preparation erases physical history but need not erase evidence about a fixed unknown parameter. Successive excursions are then conditionally Markov-renewal, rather than independent cycles. They may lead to several closed boundary classes with different durations. Their physical-time gains are ratios formed separately in each class and only then mixed by absorption probabilities. Near a recurrent-rank change, the entire ergodic projection can be discontinuous even though the prediction or decision under study remains uniformly well behaved.

The main result, Theorem~\ref{thm:transfer}, gives an exact criterion for this task-sensitive regularity. It does not postulate a spectral gap or a uniform bound for a group inverse. Under explicit compactness, finite-response continuity and return-uniform-integrability assumptions, it proves the equivalence of gain continuity, a compatibility condition at every limiting ergodic projection, uniform finite-acquisition convergence, uniform normalized-discount convergence, and a uniform approximate-corrector property. Thus it distinguishes a genuine average-risk discontinuity from a dynamically slow direction whose task amplitude vanishes. The criterion applies to all entering labels; a statement about one fortuitous initial distribution is weaker.

The quantitative part is based on the complete conditional excursion. If $P$ is its exit-label matrix, $T_{ij}=\E_i[\tau\one_{\{I_1=j\}}]$ its duration--exit matrix, $r$ its reward vector and $g$ its physical-time gain, the relevant centered reward is $r-Tg$. An approximate solution of
\begin{equation}\label{eq:intro-poisson}
 (I-P)u=r-Tg
\end{equation}
controls the discrepancy at a prescribed number of real physical calls. Duration and exit label may be correlated. This formulation avoids first paying a worst-case time to a recurrent root. Equation~\eqref{eq:intro-poisson} belongs to classical multichain semi-Markov bias theory; the present use is a uniform approximate-corrector criterion, its converse through recurrent-rank degeneration, and its coupling to one retained causal program and actual predictive occupation.

The result has three consequences relevant to the foundational problem. First, on an entire compact legal program class satisfying the criterion, a robust fixed-memory optimum is attained and is determined by finite model subfamilies. No recurrence constraint is inserted into that conclusion. Second, on finite algebraic instruments, task compatibility is decidable from the primitive model and yields an algebraic finite-acquisition modulus. Even without compatibility, a separate semialgebraic selection theorem supplies a near-optimal program whose recurrence cost grows at most polynomially in the target error, under a stated common-support hypothesis. Third, applying the criterion to bounded predictive tests transfers the actual occupation law. It can then be combined with a small-ball witness, a global candidate-valid cover and an actual-law update estimate to obtain matching causal and checkpoint resource curves.

The sharp examples are not substitutes for this theorem. They test it. In a rare-revelation decision experiment with revelation probability $|\theta|^k$, two retained labels are sufficient for vanishing worst-model risk, while one label is not. The exact finite-acquisition value is of order $N^{-1/k}$. The group inverse and the exact task bias are unbounded, but the approximate-corrector price has order $a^{-(k-1)}$, giving precisely this rate. A fixed noncommuting quantum instrument realizes the same task mechanism through a change of preparation rank. A different raw realization uses action-uniform drift for unbounded queue excursions. Neither realization is used in proving the general theorem.

For a single four-call scored protocol, the diagnostic task can be combined with a fresh geometric mark and a fixed inaccessible calibration sign. If the actually loaded mark has quantization error $q_M$, the optimized retained-memory value is bounded below and above by
\begin{equation}\label{eq:intro-law}
 q_M+\delta^2+D_n(e)
 \quad\hbox{and}\quad
 q_{\lfloor M/2\rfloor}+\delta^2+D_n(e),\qquad N=4n.
\end{equation}
Here $D_n(e)$ is an exact all-strategy diagnostic value and $e$ is a real reduction in revelation intensity, not an allowed tolerance silently promoted to a lower bound. Its common-task causal deficiency is comparable to $e$. Under acquired mass and cover bounds this gives
\[
 \calR(N,M,\delta,\varepsilon)
 \asymp M^{-2/d}+N^{-1/k}+\delta^2+\varepsilon^{1/k}.
\]
The fractional deficiency exponent is an attained task-sensitive amplification. General defect transfer, in contrast, requires regularity at both endpoints. The distinction is maintained throughout.

\subsection*{Relation to earlier theory}
The linear algebra of multichain evaluation is classical. Denardo and Fox~\cite{DenardoFox} treat multichain Markov-renewal programs; Federgrun and Schweitzer~\cite{FedergrunSchweitzer} study undiscounted Markov-renewal fixed points; Lamond and Puterman~\cite{LamondPuterman} develop generalized-inverse methods for Markov decision processes. Group inverses, classwise policy evaluation and Poisson equations are not new contributions here. In particular, replacing a correlated holding-time mark by its conditional mean independently of the exit state is not a permissible simplification of that theory.

The distinction in Theorem~\ref{thm:transfer} is between cost-universal recurrence and a specified continuous family of tasks. It gives a converse and a quantitative bridge on a compact family of raw causal responses, at an unchanged observer-state budget. Robust compact intersection is then a consequence, not a minimax interchange. Semialgebraic value asymptotics have important predecessors in Bewley and Kohlberg~\cite{BewleyKohlberg,BewleyKohlbergStationary}; real-algebraic choice and growth are standard~\cite{BasuPollackRoy}. We use them to certify task-compatible retained observers and to quantify a separate recurrence-exhaustion schedule. No efficient global synthesis algorithm is claimed.

Minimal predictive states and comparison of experiments have a long history~\cite{ShaliziCrutchfield,Blackwell,Fritz}. Our quotient lemma states its Borel realization assumptions rather than deriving them from positivity. Quantization and finite-memory learning are also established subjects~\cite{GrafLuschgy,HellmanCover}. The new resource statement identifies the actual occupation measure and the task-visible temporal price before applying those tools. The primitive instrument constructions below are deliberately explicit so that none of these layers can replace another by notation alone.

\section{Experiments, programs and predictive content}
\subsection{Raw experiments and information pattern}
All physical, action, report and history spaces are standard Borel. A model parameter $\theta\in\Theta$ is fixed for the entire execution and is not an argument of the observer's program. A prepared causal experiment consists of a preparation law and normalized nonnegative marked kernels
\[
 K_\theta(dx',dy,de,d\ell\mid x,a).
\]
Here $e$ includes visible protocol flags, $\ell\in[0,1]$ is a performed score, and every transition consumes one physical call. More general positive integer costs can be expanded into marked calls when the expansion is part of the given interface. ``Positive'' means positivity preserving and normalized, not that every report has strictly positive probability. The scoring wire need not be feedback; its visibility is specified as part of the experiment. In our diagnostic examples it is never supplied to the controller or simulator. Strategies may use only prior visible reports, their current retained label, and fresh independent randomization.

An autonomous $m$-label program, $m\le M$, comprises an initial label law, action rows, report-dependent update rows and task readouts. One tuple is selected before $\theta$ is chosen. A fresh random draw used to sample a row is not retained unless encoded in the next label. There is no extra seed, history tape, exact posterior register or free internal clock. Visible protocol phases supplied by the physical experiment are distinguished from an internally stored phase. Readout numbers belong to a specified compact set; their bit representation is a separate resource.

Physical preparations delimit excursions. Conditional on $\theta$ and the entering retained label $i$, the complete future excursion law is independent of the earlier physical history. Usually the physical preparation law is independent of $i$ and only the observer carries $i$. A prescribed label-dependent preparation is also permitted if explicitly included among the primitive legal preparation rows. This assumption is stronger than merely selecting an arbitrary return time in a hidden process. It implies that the complete conditional excursion law
\[
 \calK_x(i;d\tau,dR,dj,d\mathsf w),\qquad x=(\theta,\gamma),
\]
depends only on $i$ and $x$. Here $\tau\ge1$ is integer physical duration, $R\in[0,\tau]$ is the sum of the performed scores, $j$ is the exiting retained label, and $\mathsf w$ includes actions, reports, flags and audit marks. The observer's label is not overwritten at preparation. The triples $(\tau,R,j)$ can be arbitrarily dependent.

\subsection{Compact response families}
The following conditions will be used with a compact metric space $X$ of model--program pairs, and occasionally with an additional compact test parameter.

\begin{definition}[Continuous response family]\label{def:responses}
A family has continuous finite responses if visible finite-call cylinder probabilities and the expected bounded task readouts used below are continuous in $x$, for every entering label. This does not require total-variation continuity of a deterministic real-valued audit mark as its model value changes. It has a uniform return envelope if
\begin{equation}\label{eq:UI}
 \lim_{L\to\infty}\sup_{x,i}\E_{x,i}[\tau\one_{\{\tau>L\}}]=0.
\end{equation}
Write $\bar\mu=\sup_{x,i}\E_{x,i}\tau<\infty$ and
\begin{equation}\label{eq:tail-def}
 B_N=\sup_x\frac1N\sum_{t=1}^N\max_i\E_{x,i}(\tau-t)_+.
\end{equation}
All entering rows, including rows unreachable for a particular program, occur in these suprema.
\end{definition}
Condition~\eqref{eq:UI} is not an on-policy assumption. Altering an unreachable row can change the declared program class; it is not a free normalization. It implies $B_N\to0$ by a uniform tail split and Ces\`aro averaging. A useful raw sufficient condition is given next.

\begin{lemma}[Reused-row finite-response continuity]\label{lem:response}
Suppose the action and retained-label sets are finite. For each finite intervention word, the physical response has an $L^1$ density relative to fixed finite product reference measures on its distinct report coordinates, continuous in the model parameter in $L^1$. Let the program's finitely many simplex-valued report rows vary weak-star in the corresponding $L^\infty$ spaces, and let its finite action rows and compact readouts vary continuously. Assume the underlying $L^1$ spaces are separable. Then the legal closed program space is compact metrizable and finite responses are continuous. If~\eqref{eq:UI} holds, the excursion arrays $P,T,r$ defined below are continuous as well.
\end{lemma}
\begin{proof}
Positivity, row sums and any declared closed linear row restrictions are weak-star closed. Banach--Alaoglu and separability give compact metrizability of each row ball and hence of the finite product. A finite response is a finite sum of integrals of an $L^1$ density against products of bounded row entries on distinct coordinates. Approximate the density in $L^1$ by finite linear combinations of rectangle indicators. The resulting integrals are products of weak-star continuous scalar pairings. The approximation error is bounded uniformly by the $L^1$ error, even when exactly the same report row is reused at several times. There is no multiplication of weak-star limits on a single coordinate. Model continuity follows by the same uniform bound. Truncating an excursion at $L$ makes its exit probabilities, duration marks and rewards finite-response expressions. The discarded duration and reward are bounded by a constant times $\E[\tau\one_{\{\tau>L\}}]$, uniformly by~\eqref{eq:UI}. Passing $L\to\infty$ proves continuity of the arrays.
\end{proof}
Marginal domination alone is insufficient for this lemma. Finite report spaces satisfy the product condition automatically when their raw probabilities are continuous. Nondominated direct constructions will be identified separately and are not consequences of this compactness argument.

\subsection{Executable tests and the quotient}
A future test specifies a legal continuation strategy, a stopping rule within its stated budget and a bounded readout of permitted future reports or a designated audit. Histories with different available actions, protocol phases or remaining resources cannot be equated merely because some numerical predictions agree. Fix a determining countable collection $(f_j)$ of such tests, including the required interface coordinates, and their conditional predictions $q_j(h)$. Conditional versions and admissible histories are fixed once; no claim is made about arbitrary null histories under unspecified versions.

\begin{lemma}[Realized predictive quotient]\label{lem:quotient}
Suppose $q(h)=(q_j(h))_j$ is Borel, its image $S$ is standard Borel, and it admits a Borel section $s:S\to H$. Suppose equality of these coordinates determines every designated future test, including one-step continuations. Then equality of $q$ is exactly predictive equivalence for this test family. The quotient is realized by $S$ and is minimal among Borel statistics through which all the coordinates factor. If the report and continuation kernels are Borel, jointly in the action, and constant on these fibers, they descend to Borel kernels on $S$.
\end{lemma}
\begin{proof}
The determining assumption gives both implications in the equivalence. If $Z$ is another statistic and $q_j=\widetilde q_j\circ Z$ with Borel coordinate factors, then $q=\widetilde q\circ Z$ into the countable product; its restriction to $S$ gives the factorization on the realized range. Conversely any statistic factoring through $q$ cannot distinguish predictive-equivalent histories. For descended kernels evaluate the original kernel at $s(p)$. The section makes this evaluation Borel, including the action variable, and fiber invariance makes it independent of the section. The same argument applies to test evaluations and updates.
\end{proof}
Image regularity and a section are hypotheses, not consequences of positivity or of Lemma~\ref{lem:response}. A fixed-prior predictive quotient is not automatically a minimal statistic uniformly over an unknown parameter family. The resource theorems use a common predictive realization only when its model-independent initialization and updates have been verified.

\subsection{Resources and morphisms}
A causal simulator is parameter independent and includes a strategy lift, report transformation, internal state, preparation and clock map, compatible actions and a task readout. Its output must be due before each target decision that uses it. A simulator that stores $s$ labels together with an $m$-label observer uses at most $sm$ labels, with any protocol clock included in one of these coordinates. Program description, workspace, acquisition count and physical time do not follow from this label bound. A common-task deficiency compares the complete scored interface with an unchanged environment-owned audit wire; it is not unrestricted replacement of the ground-truth task. Precise finite-call and long-run error statements are given in Section~\ref{sec:composition}.

\section{Task-visible causal transfer}\label{sec:transfer}
Fix a continuous response family on $X$ as in Definition~\ref{def:responses}, with $m$ labels. Define the continuous excursion arrays
\begin{equation}\label{eq:arrays}
 P_{ij}=\Pp_i(I_1=j),\quad T_{ij}=\E_i[\tau\one_{\{I_1=j\}}],\quad
 d=T\one,\quad D=\diag(d),\quad r_i=\E_iR.
\end{equation}
Matrices act on column rewards. Put
\begin{equation}\label{eq:time-change}
 S=I-D^{-1}(I-P),\quad c=D^{-1}r,\quad
 \Pi_A=\lim_{n\to\infty}\frac1n\sum_{j=0}^{n-1}A^j,
 \qquad g=\Pi_Sc.
\end{equation}
The stochasticity of $S$ follows from $d_i\ge1$: its diagonal is $1-(1-P_{ii})/d_i\ge0$. Every finite stochastic matrix has a semisimple eigenvalue at one and a Ces\`aro projection, regardless of periodicity. Let $J_N(x,i)$ be the expected sum of the first $N$ physical-call scores divided by $N$, from boundary label $i$, and let
\[
 J_\beta(x,i)=(1-\beta)\E_{x,i}\sum_{t\ge1}\beta^{t-1}\ell_t,
 \qquad 0<\beta<1.
\]
For $a>0$ define the task-corrector price
\begin{equation}\label{eq:corrector-price}
 \calC_X(a)=\sup_{x\in X}\inf\left\{\osc(u):
  \norm{r-Tg-(I-P)u}_\infty\le a\right\}.
\end{equation}
The infimum is $+\infty$ if the constraint is empty; we will show it is never empty pointwise. Oscillation is $\max_i u_i-\min_i u_i$. The exact price at $a=0$ need not be uniformly finite.

\begin{theorem}[Task-visible causal transfer]\label{thm:transfer}
Under the hypotheses above, $J_N(x)\to g(x)$ and $J_\beta(x)\to g(x)$ pointwise. The following five statements are equivalent:
\begin{enumerate}
\item[(i)] The vector $g:X\to\R^m$ is continuous.
\item[(ii)] Whenever $x_j\to x$ and $\Pi_{S(x_j)}\to E$, one has $Eg(x)=g(x)$.
\item[(iii)] $\sup_x\norm{J_N(x)-g(x)}_\infty\to0$.
\item[(iv)] $\sup_x\norm{J_\beta(x)-g(x)}_\infty\to0$ as $\beta\uparrow1$.
\item[(v)] $\calC_X(a)<\infty$ for every $a>0$.
\end{enumerate}
For every $a>0$, without assuming (i), the following inequalities hold whenever their right sides are finite:
\begin{align}
 \sup_x\norm{J_N-g}_\infty
 &\le a+\frac{\calC_X(a)}{N}+B_N,\label{eq:physical-bound}\\
 \sup_x\norm{J_\beta-g}_\infty
 &\le a+(1-\beta)\calC_X(a)
 +(1-\beta)^2\sum_{N\ge1}N\beta^{N-1}B_N.\label{eq:discount-bound}
\end{align}
If $\sup_{x,i}\E_i\tau^{1+p}\le K$ for $0<p\le1$, the last tail terms are at most $mKN^{-p}$ and $mK(1-\beta)^p$, respectively. The factor $m$ can be omitted when one common stochastically dominating lifetime has this moment bound. All comparisons use the same program and the same retained-label budget.
\end{theorem}
Condition (ii) permits a change of recurrent rank. It asks only that each limiting projection fix the limiting \emph{task gain}; it does not require the projection itself to converge. Uniformity is asserted on the entire declared compact family $X$, not on arbitrary experiments lacking the response or return assumptions.

\begin{lemma}[Marked bias and physical-time gain]\label{lem:marked}
The vector $g$ is $P$-harmonic. If $C$ runs over the closed classes of $P$, $\pi_C$ is its stationary row on $C$, and $w_C(i)$ is the absorption probability from $i$, then
\begin{equation}\label{eq:class-gain}
 g_i=\sum_C w_C(i)\frac{\pi_Cr}{\pi_Cd}.
\end{equation}
Moreover $b=r-Tg$ satisfies $\Pi_Pb=0$, and $(I-P)u=b$ has a finite solution.
\end{lemma}
\begin{proof}
$S$ has the same off-diagonal graph, closed classes and absorption probabilities as $P$. On $C$ its stationary row is $\pi_C D/(\pi_Cd)$, which proves~\eqref{eq:class-gain} and harmonicity. On a closed class $g_j=g_C$ for every possible exit, so $(Tg)_i=d_i g_C$ there and $\pi_Cb=0$. Every row of $\Pi_P$ is an absorption mixture of such rows, hence $\Pi_Pb=0$. The range of $I-P$ is the kernel of $\Pi_P$. For example the group inverse
\[
 G_P=(I-P+\Pi_P)^{-1}-\Pi_P
\]
satisfies $(I-P)G_P=I-\Pi_P$, so $u=G_Pb$ is a solution. These are finite-state multichain identities, not a contraction or irreducibility assumption.
\end{proof}

\begin{lemma}[Approximate marked corrector at physical time]\label{lem:physical}
For a fixed experiment--program pair and any $u$ with $\norm{b-(I-P)u}_\infty\le a$,
\[
 \norm{J_N-g}_\infty
 \le a+\frac{\osc(u)}N+\frac1N\sum_{t=1}^N\max_i\E_i(\tau-t)_+.
\]
\end{lemma}
\begin{proof}
Let $I_k$ be the entering label of excursion $k$, $\tau_k$ its duration, and $R_k$ its reward. Set $S_0=0$ and $S_k=\sum_{j<k}\tau_j$. The boundary filtration $\mathcal F_k$ contains precisely the past before excursion $k$ is run. Define
\[
 K_N=\min\{k\ge1:S_k\ge N\}.
\]
Then $K_N\le N$ and $\{k<K_N\}=\{S_k<N\}\in\mathcal F_k$. With $e=b-(I-P)u$, the conditional mean of
\[
 R_k-\tau_k g(I_{k+1})+u(I_{k+1})-u(I_k)
\]
is $e(I_k)$. Summing with these predictable indicators is legitimate: the number of terms is bounded by $N$ and their absolute expectations are finite.

Since $Pg=g$, $g(I_k)$ is a bounded martingale. Summation by parts gives
\[
 \sum_{k<K_N}\tau_k g(I_{k+1})
 =S_{K_N}g(I_{K_N})-
   \sum_{k<K_N}S_k\bigl(g(I_{k+1})-g(I_k)\bigr).
\]
The second sum has expectation zero because its coefficients are predictable and less than $N$. Also $\E g(I_{K_N})=g(I_0)$. Consequently the expected reward of the completed crossing excursion, less $Ng(I_0)$, is the sum of a residual bounded by $aN$, a corrector boundary term bounded by $\osc(u)$, and $\E[(S_{K_N}-N)g(I_{K_N})]$.

The reward after call $N$ but before the crossing excursion ends lies in $[0,S_{K_N}-N]$. So does $(S_{K_N}-N)g(I_{K_N})$. Their difference has absolute expectation at most the expected overshoot, not twice that quantity. There is at most one boundary at any integer physical time $s$. Conditioning there gives
\[
 \E(S_{K_N}-N)\le
 \sum_{s=0}^{N-1}\max_i\E_i(\tau-(N-s))_+.
\]
Division by $N$ proves the assertion. The proof allows dependence among reward, exit label and duration; it never replaces $T$ by $DP$.
\end{proof}

\begin{proof}[Proof of Theorem~\ref{thm:transfer}]
Pointwise finite-acquisition convergence follows from Lemma~\ref{lem:physical} with the exact finite corrector of Lemma~\ref{lem:marked} and $a=0$. Its tail term vanishes by integrability. Taking the infimum over correctors and then the supremum proves~\eqref{eq:physical-bound}; an arbitrarily small excess over the infimum can be removed. For bounded scores Abel summation yields the exact identity
\begin{equation}\label{eq:abel}
 J_\beta=(1-\beta)^2\sum_{N\ge1}N\beta^{N-1}J_N.
\end{equation}
The weights sum to one, proving pointwise discount convergence and~\eqref{eq:discount-bound}.

For the moment statement,
\[
 \sum_{t=1}^N\max_i\E_i(\tau-t)_+
 \le\sum_i\E_i[\tau\min(\tau,N)]\le mK N^{1-p}.
\]
A common dominating lifetime replaces the sum over labels by one expectation. If $G$ is geometric on $\N$ with success probability $1-\beta$, Jensen's inequality gives
\[
 (1-\beta)^2\sum_{N\ge1}N^{1-p}\beta^{N-1}
 =(1-\beta)\E G^{1-p}\le(1-\beta)^p.
\]

Finite-response continuity makes each $J_N$ continuous; a uniformly convergent bounded discounted series makes $J_\beta$ continuous. Thus (iii) or (iv) implies (i). Equations~\eqref{eq:physical-bound}--\eqref{eq:discount-bound}, first with a fixed small $a$ and then with large $N$ or $\beta$ close to one, give (v)$\Rightarrow$(iii),(iv).

Suppose (i) holds. Then $b(x)=r(x)-T(x)g(x)$ is continuous and bounded, say by $B$. By Lemma~\ref{lem:marked},
\[
 A_n(x)b(x):=\frac1n\sum_{k=1}^nP(x)^k b(x)\longrightarrow0
\]
pointwise. This convergence is uniform. Indeed, for each $x_0$ choose a block length $q$ with $\norm{A_q(x_0)b(x_0)}<\epsilon$. The same bound, with $2\epsilon$, holds in a neighborhood by finite-dimensional continuity. Decomposing any $n$ into blocks of length $q$ and a remainder, and using stochastic contraction, gives $\norm{A_n b}\le2\epsilon+qB/n$ there. A finite cover of $X$ supplies a uniform bound on the chosen lengths. Now put
\[
 u_n(x)=\sum_{k=0}^{n-1}(1-k/n)P(x)^k b(x).
\]
Direct telescoping gives $(I-P)u_n=b-A_n b$ and $\osc(u_n)\le(n+1)B$. Choosing one sufficiently large $n$ for a given $a$ proves (v). No uniform bound on the exact group inverse has been used.

Finally, if $x_j\to x$ and $\Pi_{S(x_j)}\to E$, then $ES(x)=S(x)E=E$. Hence $E\Pi_{S(x)}=E$. Continuity of $c$ gives $g(x_j)\to Ec(x)=Eg(x)$ along that subsequence. Thus (i) implies (ii). Conversely the projections belong to a compact set of stochastic matrices, so every subsequence has such a further subsequence. Condition (ii) forces all resulting limits of $g(x_j)$ to equal $g(x)$, which proves (i).
\end{proof}

\begin{corollary}[Cost-universal sublevels as a special case]\label{cor:group}
If $\sup_x\norm{G_{P(x)}}_\infty\le H$, then
\[
 \sup_x\norm{J_N-g}_\infty\le\frac{2\bar\mu H}{N}+B_N.
\]
On a compact family of stochastic matrices, uniform boundedness of $G_P$, continuity of $\Pi_P$, locally constant recurrent rank, and uniform Ces\`aro convergence for all bounded rewards are equivalent. None is necessary for the single-task theorem.
\end{corollary}
\begin{proof}
Each coordinate of $r$ and $Tg$ lies in $[0,d_i]$, so $\norm b_\infty\le\bar\mu$. The exact solution $G_Pb$ has oscillation at most $2H\bar\mu$. For the classical equivalence use
\[
 \sum_{k=0}^{n-1}(P^k-\Pi_P)=(I-P^n)G_P.
\]
Uniform convergence implies continuity of the projection and hence locally constant integer trace. Fixed nullity at one permits a locally fixed spectral contour: otherwise an additional eigenvalue could approach one and increase nullity at the limiting stochastic matrix, whose eigenvalue one is semisimple. The contour projection and $G_P$ are then continuous; a finite cover bounds the norm. A bounded sequence of group inverses is precompact in finite dimension, and passing their defining polynomial identities to a limit shows that a norm sublevel is closed. These standard facts are recorded to delimit, not enlarge, the scope of the task theorem~\cite{Meyer,LamondPuterman}.
\end{proof}

\begin{corollary}[Robust attainment without an imposed recurrence bound]\label{cor:robust}
Let $\Theta$ and a nonempty legal program class $\Gamma_M$ be compact metric spaces, with continuous initial-label law included in the program. Assume Theorem~\ref{thm:transfer}(i) on $X=\Theta\times\Gamma_M$. Then
\begin{equation}\label{eq:robust-value}
 V_M=\min_{\gamma\in\Gamma_M}\sup_{\theta\in\Theta}g(\theta,\gamma)
 =\sup_{\substack{F\subset\Theta\\0<|F|<\infty}}
       \min_{\gamma\in\Gamma_M}\max_{\theta\in F}g(\theta,\gamma),
\end{equation}
where scalar gains use that initial law. The optimal expected $N$-call and normalized-discount values converge to $V_M$ with the errors in~\eqref{eq:physical-bound}--\eqref{eq:discount-bound}. Cluster points of asymptotically optimal programs are average optimal.
\end{corollary}
\begin{proof}
Joint gain continuity makes the worst-model objective continuous on the compact program space. For a proposed level $v$, the closed sets $\{\gamma:g(\theta,\gamma)\le v\}$ have the finite-intersection property precisely when every finite-model problem admits that level. Compactness proves~\eqref{eq:robust-value}; there is no exchange of minimization with maximization. Uniform same-program errors pass through both extrema. Compactness and continuity then prove the assertion about cluster points.
\end{proof}
When $\Gamma_M$ is the full legal class, this is an unrestricted result within the stated raw interface. When it is an architecture, the conclusion is only about that architecture. The task criterion must be proved on the whole class used in the optimization.

\section{Primitive certification and operational recurrence exhaustion}\label{sec:algebraic}
The exact criterion in Theorem~\ref{thm:transfer} does not require finite-dimensional physical dynamics. An effective consequence does. In this section all parameter and program sets are compact semialgebraic sets defined over the real algebraic numbers, all alphabets are finite, and primitive marked-response matrix entries and task readouts are continuous semialgebraic functions. Thus the finite-response continuity hypotheses remain in force. They may be rational after introducing algebraic auxiliary variables. A uniform primitive termination certificate is part of the input. This is not an algorithm for arbitrary computable Borel kernels.

\begin{lemma}[Finite instrumental resolvents]\label{lem:resolvent}
For a finite classical physical experiment with retained labels, let $A$ be the substochastic alive matrix on physical-state--label pairs, $B$ the one-step exit matrix to boundary labels, and $E$ the preparation matrix. Assume $\rho(A)<1$ uniformly. If $v$ is the expected one-step score from an alive pair, then
\begin{equation}\label{eq:resolvents}
 P=E(I-A)^{-1}B,\qquad
 T=E(I-A)^{-2}B,\qquad
 r=E(I-A)^{-1}v.
\end{equation}
The same formulas hold for finite-dimensional quantum instruments after replacing $A$ by the completely positive alive superoperator on the direct sum of Hermitian physical blocks indexed by observer labels, and replacing matrix rows by preparation states and exit/score functionals. In either case the excursion arrays are semialgebraic functions of the primitive data.
\end{lemma}
\begin{proof}
An excursion exiting after $j\ge1$ calls contributes $EA^{j-1}B$ to $P$ and $jEA^{j-1}B$ to $T$. The absolutely convergent sums are $(I-A)^{-1}$ and $(I-A)^{-2}$; the reward sum is the alive resolvent applied to the one-step score. In the quantum case work with unnormalized states. Composition and summation of marked completely positive maps are linear on the finite real vector space of Hermitian blocks. An exit functional applied to $A^{j-1}$ has exactly the same duration weight $j$. A uniform trace-decay certificate gives absolute convergence. Inversion of a finite nonsingular real matrix is rational in its entries. No normalization at a zero-probability report or closure of an unbounded operator is involved.
\end{proof}
A fixed preparation call can be included as an extra transient block; it must not be removed from $T$ or $r$. A positive classical or quantum Lyapunov functional giving $\norm{A^j}\le C\alpha^j$, $\alpha<1$, is a sufficient uniform termination input and supplies every finite return moment.

\begin{theorem}[Decidable task regularity and algebraic modulus]\label{thm:algebraic}
For the finite algebraic instruments just described, gain continuity on a specified compact semialgebraic family $X$ is decidable. If it holds, there exist effectively obtainable constants $a_0>0$, $C<\infty$ and rational $\kappa\ge0$ such that
\begin{equation}\label{eq:algebraic-price}
 \calC_X(a)\le C a^{-\kappa},\qquad 0<a<a_0.
\end{equation}
Consequently the acquisition error is
\begin{equation}\label{eq:algebraic-rate}
 O\bigl(N^{-1/(\kappa+1)}\bigr)+B_N,
\end{equation}
and the discount error is $O((1-\beta)^{1/(\kappa+1)})$ plus its tail term. The optimal robust average value and whether it is attained on the declared semialgebraic program class are decidable even when gain continuity fails. These conclusions are decision and global certification results, not polynomial-time synthesis claims.
\end{theorem}
\begin{proof}
By Lemma~\ref{lem:resolvent}, $P,T,r,S,c$ are semialgebraic. The ergodic projection is semialgebraic on each rank stratum. One direct formula is the coefficient of $z^{-1}$ in $(zI+I-S)^{-1}$: if the determinant has first nonzero term $a_qz^q$, semisimplicity at zero of $I-S$ makes that coefficient the coefficient of $z^{q-1}$ in the adjugate divided by $a_q$. There are only finitely many possible $q$. Thus the graph of $g=\Pi_Sc$ is a finite semialgebraic union. Continuity can be expressed by the first-order condition
\[
 \forall x\in X\ \forall\epsilon>0\ \exists\eta>0\
 \forall y\in X:\quad \norm{x-y}<\eta\Longrightarrow
 \norm{g(x)-g(y)}<\epsilon.
\]
Quantifier elimination over real closed fields decides it.

For the price, impose $u_1=0$, introduce upper and lower bounds for its coordinates, and encode the residual inequality in~\eqref{eq:corrector-price}. The infimum of the resulting finite linear-inequality problem and its supremum over $X$ have semialgebraic graphs, by quantifier elimination; an infimum need not be selected to make this statement. Under continuity they are finite for every $a>0$ by Theorem~\ref{thm:transfer}. A finite one-variable semialgebraic function has polynomial growth at an endpoint. The semialgebraic growth theorem, together with effective real-algebraic decomposition, supplies~\eqref{eq:algebraic-price} with rational exponent and algebraic bounds~\cite{BasuPollackRoy}. Balancing $a$ with $Ca^{-\kappa}/N$ proves~\eqref{eq:algebraic-rate}; the same balance works for discount.

Finally, the graph of the worst-model gain is semialgebraic using quantified upper bounds and arbitrarily close lower witnesses. Its infimum on the program set is therefore a real algebraic number when all defining coefficients are algebraic. The sentence asserting a program with that value decides attainment. The argument applies to a discontinuous objective as well, but then asserts neither compact attainment nor uniform finite-acquisition convergence unless the corresponding criterion holds. The algebraic assertions use standard real-algebraic machinery, not a numerical sample of the parameter continuum.
\end{proof}
The exponent in~\eqref{eq:algebraic-price} is a certified possible exponent, not necessarily the sharp exponent. The example in Section~\ref{sec:sharpness} identifies the exact price order directly. Algebraic descriptions of row probabilities are finite programs for real-arithmetic sampling; conversion to a particular bit-level random sampler, its variable running time and its approximation errors are separate implementation obligations.

\subsection{An actionable exhaustion beyond compact attainment}
A different conclusion remains available when the average objective is discontinuous. It must not be confused with Theorem~\ref{thm:algebraic}'s compatible case.

\begin{theorem}[Polynomial near-optimal recurrence schedule]\label{thm:schedule}
Assume the finite algebraic setting above, a uniform bound $\bar\mu$ on mean physical duration and a uniform return envelope over all legal programs. Suppose also that for each fixed program $\gamma$, the boundary support graph of $P(\theta,\gamma)$ is independent of $\theta\in\Theta$. Let
\[
 V=\inf_{\gamma\in\Gamma_M}\sup_{\theta\in\Theta}g(\theta,\gamma).
\]
There are $a_0>0$, $C<\infty$, rational $\kappa\ge0$ and a semialgebraic choice of common programs $\gamma_a$, $0<a<a_0$, such that
\begin{equation}\label{eq:schedule}
 \sup_\theta g(\theta,\gamma_a)\le V+a,
 \qquad
 \sup_\theta\norm{G_{P(\theta,\gamma_a)}}_\infty\le Ca^{-\kappa}.
\end{equation}
In particular the same program $\gamma_a$ satisfies, for every $N$,
\begin{equation}\label{eq:schedule-physical}
 \sup_\theta J_N(\theta,\gamma_a)
 \le V+a+2\bar\mu Ca^{-\kappa}/N+B_N,
\end{equation}
with the analogous discount bound. The data determine whether the unrestricted average infimum is attained. Without task compatibility, no assertion is made that the unrestricted finite-$N$ optimal values converge to $V$.
\end{theorem}
\begin{proof}
For fixed $\gamma$, the common graph fixes the number of closed classes as $\theta$ varies. The spectral-contour argument in Corollary~\ref{cor:group}, applied on compact $\Theta$, makes
\[
 H(\gamma)=\max_\theta\norm{G_{P(\theta,\gamma)}}_\infty
\]
finite. Across programs this function can be unbounded and discontinuous, but it is semialgebraic by the finite rank-stratum formulas and quantifier elimination. The set of pairs $(a,\gamma)$ satisfying the first inequality of~\eqref{eq:schedule} is nonempty for every $a>0$ and semialgebraic. Semialgebraic choice supplies a curve $\gamma_a$. The finite semialgebraic function $H(\gamma_a)$ has polynomial growth as $a\downarrow0$, proving the second inequality after shrinking $a_0$. Corollary~\ref{cor:group} gives~\eqref{eq:schedule-physical}. Attainment is the first-order statement proved decidable in Theorem~\ref{thm:algebraic}.
\end{proof}
This theorem converts recurrence exhaustion into a primitive-data, near-optimal implementation schedule. It permits the required recurrence level to diverge and does not manufacture an optimizer at the limit. The common-support assumption is testable in the finite algebraic setting; it is not a consequence of continuity alone. In particular, the rank-changing rare-herald family below uses Theorem~\ref{thm:transfer}, not this common-support theorem. Semialgebraic strategy asymptotics in stochastic games have classical precedents~\cite{BewleyKohlberg,BewleyKohlbergStationary}; the schedule here applies them to one fixed-memory robust causal program and its separately charged physical-return bridge.

\section{Actually acquired occupation and causal resolution}\label{sec:geometry}
The task criterion is useful only if the relevant geometry is itself acquired by the experiment. Fix a declared exploration and a common predictive realization when one is available. At designated scored calls let $p_t=\E[Y_t\mid H_t]$ be the conditional mean of a bounded finite-dimensional target, with $p_t\in K$ for a compact $K\subset\R^q$. In a robust family this formula is interpreted modelwise, and a common executable representation must be verified separately. Define the excursion occupation measures
\[
 \nu_i(A)=\E_i\sum_{t\text{ scored in excursion}}\one_{\{p_t\in A\}}.
\]
Preparation or other zero-score calls are omitted from this numerator but included in the duration $d_i$. The actual physical-time measure from initial label $i$ is
\begin{equation}\label{eq:occupation}
 \nu_\infty^i=\sum_C w_C(i)\frac{\sum_{j\in C}\pi_C(j)\nu_j}{\pi_Cd}.
\end{equation}
All quantization identities below concern a specified prepared law. In a robust family, modelwise quantization gives a lower bound, whereas an implementable upper must use one common codebook and recursion. The measure in~\eqref{eq:occupation} is a subprobability measure. Transient boundary occupation has zero asymptotic mass. Let $\nu_N^i$ be the expected empirical scored occupation divided by $N$ physical calls, and define $\nu_\beta^i$ with normalized-discount weights. These measures are neither preparation measures nor freely chosen volume measures.

\begin{theorem}[Task-visible acquired-law criterion]\label{thm:occupation}
Assume continuous finite responses and~\eqref{eq:UI}. Let the scored feature maps be such that finite responses of every bounded Lipschitz function of $p_t$ are continuous in $x$. Then the following are equivalent: continuity of $x\mapsto\nu_\infty^i$ in the bounded-Lipschitz metric for all initial labels; uniform convergence $\nu_N^i\to\nu_\infty^i$ in that metric; uniform discounted convergence; and Theorem~\ref{thm:transfer}'s task criterion on the compact family of normalized nonnegative Lipschitz tests of $K$. Its quantitative bounds use the supremum of the corresponding approximate-corrector prices.
\end{theorem}
\begin{proof}
Use $\mathcal F=\{f:K\to[0,1]:\Lip(f)\le1\}$, which is compact in the uniform norm by Arzel\`a--Ascoli. Treat $f$ as an additional parameter and score $f(p_t)$ only at designated calls. Formula~\eqref{eq:class-gain} is exactly~\eqref{eq:occupation} tested against $f$. Finite-response continuity is joint in $(x,f)$ because uniform changes in $f$ change normalized expected scores by at most their uniform norm. Continuity of the measure implies joint continuity of the gain; conversely continuity for all $f$ on compact $X\times\mathcal F$ gives bounded-Lipschitz continuity. Apply Theorem~\ref{thm:transfer} on this product. This argument retains the true mass of every stratum and does not require a smooth coordinate chart.
\end{proof}

For a finite positive measure $\nu$ on $K$, define
\begin{equation}\label{eq:quantization}
 q_J(\nu)=\inf_{c_1,\ldots,c_J\in\conv K}
       \int\min_{1\le j\le J}\norm{p-c_j}^2\,d\nu(p).
\end{equation}
In finite-dimensional compact support the infimum is attained. The equality-of-infima formulation remains meaningful in a general Hilbert space, where attainment requires an additional argument. Put $D_K=\operatorname{diam}(K)$.

\begin{lemma}[Projection, actual mass and quantization]\label{lem:quantization}
For a fixed exploration, the best checkpoint risk with $J$ fixed response values is its conditional-variance baseline plus $q_J$ of the actual scored occupation. Every online observer with at most $J$ such values has at least this risk, even when those labels are also used for control. Further,
\begin{equation}\label{eq:q-BL}
 |q_J(\nu)-q_J(\nu')|
 \le(D_K^2+2D_K)\,d_{\rm BL}(\nu,\nu'),
\end{equation}
with a metric using bounded-plus-Lipschitz norm at most one. If $\lambda\le\nu$ has mass $w>0$ and
\[
 \lambda(B(z,r))\le C r^d\qquad(z\in\conv K,\ 0<r\le r_0),
\]
then, whenever $r=(w/(2CJ))^{1/d}\le r_0$,
\begin{equation}\label{eq:smallball}
 q_J(\nu)\ge\frac w2\left(\frac{w}{2CJ}\right)^{2/d}.
\end{equation}
A global cover of $K$ by $J$ radius-$h_J$ balls gives $q_J(\nu)\le\nu(K)h_J^2$.
\end{lemma}
\begin{proof}
Conditional squared projection gives
$\E\norm{Y-c}^2=\E\norm{Y-p}^2+\E\norm{p-c}^2$ for each measurable response, with independent randomized choices included by conditioning. For a fixed codebook the pointwise nearest response is optimal. A checkpoint can make this assignment measurably; any online response is bounded below by the same pointwise distance. The function $\min_j\norm{p-c_j}^2$ has bound $D_K^2$ and Lipschitz constant at most $2D_K$. Taking infima proves~\eqref{eq:q-BL}. The union of $J$ radius-$r$ balls has $\lambda$-mass at most $JCr^d=w/2$; the remaining mass contributes at least $wr^2/2$. The cover bound is immediate. No mass is normalized away.
\end{proof}
The baseline may vary with the exploration. The pointwise lower remains valid when the observer jointly controls and predicts, but it is then the occupation law generated by \emph{that} observer that enters.

\subsection{A counted causal implementation}
A geometric cover alone is not an executable predictor. The following assumptions specify the missing causal content. Fix an exploration using $Q$ retained control labels. The predictive realization may use a state carrier larger than $K$: write its state as $s$, and require that the conditional task mean is $p=f(s)\in K$. Suppose it has a common known initialization at physical preparation, a jointly Borel update $F(s,a,y)$ and a task map $f$ with Lipschitz constant $L_f$. When the carrier is the conditional-mean space itself, take $f$ to be the identity. All quantization measures below are measures of $p=f(s)$, not of an unrelated carrier coordinate. A budget-independent relation between exact and candidate states is specified. For every $J$ there is a globally valid $J$-center projection of error at most $h_J$ on the full candidate domain; the update, projection and numerical perturbations preserve that relation. The report kernel is continuous on that domain in the declared task or total-variation modulus. Continuity is not a substitute for the following same-report stability estimate.

Along the \emph{actual} reports and exploration actions, suppose every admissible candidate recursion satisfies
\begin{equation}\label{eq:gain-recursion}
 e_t\le L_t e_{t-1}+h_J+\delta,
 \qquad
 G_t=1+L_tG_{t-1},\quad G_{\rm prep}=1,
 \qquad e_{\rm prep}\le h_J+\delta.
\end{equation}
The predictor, but not the retained control label, is initialized at preparation. Assume
\begin{equation}\label{eq:occupied-stability}
 \sup_{x,N}\frac1N\E_x\sum_{t\le N,\ t\text{ scored}}G_t^2\le S_0<\infty.
\end{equation}
A limiting or horizon-dependent version gives the corresponding limiting or horizon-dependent bound. This is an actual-acquisition condition, not stability on a formal reachable set.

\begin{theorem}[Acquired geometry to causal memory]\label{thm:geometry}
Under these assumptions, a $QJ$-label observer implements the exploration and a predictor whose squared excess risk at every physical horizon is at most
\begin{equation}\label{eq:geometric-upper}
 \left[L_f(h_J+\delta)\sqrt{S_0}+\xi\right]^2,
\end{equation}
where $\xi$ bounds the numerical task readout error. Its retained control coordinate survives preparation. The checkpoint excess is exactly $q_J(\nu_N)$; every online $M$-label observer has excess at least $q_M(\nu_N)$ for its own actual law. Under Theorem~\ref{thm:occupation}, finite-acquisition and average quantization laws differ by~\eqref{eq:q-BL} and the task-corrector modulus.

For optimized average risk, suppose every legal $M$-label observer has baseline at least $B_*$ and its own actual occupation contains the same uniform small-ball witness parameters $(w,C,d,r_0)$. Suppose a fixed $Q$-label exploration attains $B_*$ and satisfies the global cover and stability assumptions with $h_J\le C_1J^{-1/d}$. Then, for $M\ge2Q$ and exact arithmetic,
\begin{equation}\label{eq:optimized-geometry}
 cM^{-2/d}\le\inf_{\gamma\in\Gamma_M}\calR_\infty(\gamma)-B_*
 \le C_2M^{-2/d}.
\end{equation}
The lower quantifiers are over all legal observers; the upper uses the displayed exploration.
\end{theorem}
\begin{proof}
Store $(q,c)$, the exploration label and the predictive center. Actions use only $q$, so the physical exploration law is unchanged. On an actual report, update the candidate from $c$ and project it; the next control label is updated by its original rule. There are exactly $QJ$ available pairs, including any internal clock already charged to the exploration. At preparation reset only the predictor coordinate to its common approximate prior. Induction in~\eqref{eq:gain-recursion} gives $e_t\le(h_J+\delta)G_t$. The task map and numerical readout yield pointwise root error at most $L_f(h_J+\delta)G_t+\xi$. Minkowski's inequality in the actual normalized scored occupation and~\eqref{eq:occupied-stability} give~\eqref{eq:geometric-upper}; its square, not its root, is the excess risk. Lemma~\ref{lem:quantization} and Theorem~\ref{thm:occupation} prove the checkpoint and temporal assertions.

For optimized matching, apply~\eqref{eq:smallball} to the law of each competing observer, using its at most $M$ task readouts even if the labels also implement control. This gives the uniform lower above $B_*$. For the upper choose $J=\lfloor M/Q\rfloor$, which is at least $M/(2Q)$, and use the attaining exploration with the counted product implementation. These are the extra assumptions required to pass from fixed-exploration geometry to optimized control; no such passage is automatic.
\end{proof}

\begin{proposition}[Survival-weighted expansion]\label{prop:survival}
Suppose an excursion survives each further scored step with conditional probability at most $s<1$, and every candidate-valid one-step Lipschitz multiplier is at most $L$. If $\sqrt{s}L<1$, then an actual-law stability constant as in~\eqref{eq:occupied-stability} is finite, bounded by a constant depending only on $s,L$ and the fixed number of non-scored preparation calls. Individual updates may expand, $L>1$.
\end{proposition}
\begin{proof}
At depth $k$, $G_k\le\sum_{j=0}^kL^j$. Minkowski in the weighted sequence space gives
\[
 \left(\sum_{k\ge0}s^kG_k^2\right)^{1/2}
 \le\sum_{j\ge0}L^j\left(\sum_{k\ge j}s^k\right)^{1/2}
 =\frac1{\sqrt{1-s}(1-\sqrt{s}L)}.
\]
Conditioning at each actual preparation and using at most one preparation per physical call bounds normalized cumulative occupation by this excursion sum, up to the fixed convention at preparation. No conditioning on successful survival removes its probability. Expected logarithmic contraction or random-product hypotheses may be weaker, but are not asserted here.
\end{proof}
This sufficient stability mechanism and the task-visible occupation criterion address different obstructions. The former realizes a recursive predictor; the latter permits temporal rank degeneration invisible to its test family. Neither implies the other without the displayed hypotheses.

\section{A sharp retained-information experiment}\label{sec:sharpness}
Consider a fixed unknown $\theta\in[-1,1]$, put $t=|\theta|$ and $s=\operatorname{sign}(\theta)$ when $t>0$. One diagnostic call requires a decision $u\in\{-1,+1\}$ before its report. Its score is $t\one_{\{u\ne s\}}$, with score zero at $t=0$. The score is recorded on an environment-owned audit wire and is not supplied as feedback, now or later. The visible report is $s$ with probability $\lambda(t)$ and $\bot$ otherwise, independently across physical preparations conditional on $\theta$. Here $\lambda$ is continuous, $\lambda(0)=0$, and $\lambda(t)>0$ for $t>0$. The raw probabilities and scores are continuous on $[-1,1]$; the vanishing factor $t$ removes the apparent sign discontinuity at zero. The observer's label is retained.

\begin{theorem}[Exact diagnostic value and sharp corrector rate]\label{thm:diagnostic}
For every $n\ge1$ and every $M\ge2$, even allowing arbitrary full-history competing policies, the minimax $n$-call risk is
\begin{equation}\label{eq:diagnostic-value}
 D_n(\lambda)=\sup_{0\le t\le1}\frac{t}{2n}
                   \sum_{j=0}^{n-1}(1-\lambda(t))^j.
\end{equation}
One stationary two-label program attains it for all $n$: initialize its sign label fairly, retain it on $\bot$, and replace it by every conclusive report. Its normalized-discount value is
\begin{equation}\label{eq:diagnostic-discount}
 D_\beta(\lambda)=\sup_{0\le t\le1}
 \frac{t(1-\beta)}{2[1-\beta(1-\lambda(t))]}.
\end{equation}
For $\lambda(t)=t^k$ with integer $k\ge2$, these values are respectively of orders $n^{-1/k}$ and $(1-\beta)^{1/k}$. The two-label program satisfies Theorem~\ref{thm:transfer} on the entire parameter interval, although its group-inverse norm and its exact task bias are unbounded. Its approximate-corrector price satisfies
\begin{equation}\label{eq:sharp-price}
 \calC(a)\asymp_k a^{-(k-1)}\qquad(a\downarrow0).
\end{equation}
A stationary one-label program has minimax average risk $1/2$ when $\lambda(1)=1$; two retained labels have minimax average risk zero.
\end{theorem}
\begin{proof}
Fix $t>0$ and average the two models $\theta=t,-t$ with equal weights, solely for a minimax lower bound. Until the first conclusive report their observable histories have the same law. Their decisions may depend on arbitrary past actions and private fresh randomization, but the conditional average error probability before that report is $1/2$. At decision $j+1$, this event has probability $(1-\lambda(t))^j$. Therefore every policy has worst-model risk at least the expression in~\eqref{eq:diagnostic-value}, for each $t$. The two-label program has exactly this risk under each sign: the initial fair guess remains fair before revelation and is correct afterwards. It attains the lower for every $t$ simultaneously, proving the optimized claim. Summing the discounted geometric series gives~\eqref{eq:diagnostic-discount}.

For $\lambda(t)=t^k$,
\[
 \frac{t}{2n}\sum_{j<n}(1-t^k)^j
 \le\frac12\min(t,t^{1-k}/n).
\]
Splitting at $t=n^{-1/k}$ gives the upper rate. At that value of $t$, the sum divided by $n$ is bounded below by a positive absolute constant (with $n=1$ handled directly), giving the lower. The discount bounds follow by splitting at $t=(1-\beta)^{1/k}$; the denominator is $(1-\beta)+\beta t^k$, and $\beta\le1/2$ is absorbed by a fixed constant.

For positive $\theta$, order the labels as correct, incorrect. Then
\[
 P_t=\begin{pmatrix}1&0\\t^k&1-t^k\end{pmatrix},\quad
 r_t=\binom{0}{t},\quad \tau=1,\quad g_t=0.
\]
For negative $\theta$ exchange the labels; at zero $P_0=I$ and $r_0=0$. Thus the gain is identically zero from either entering label while the recurrent rank changes. The group inverse has norm of order $t^{-k}$ and an exact corrector needs oscillation $t^{1-k}$. An approximate corrector with residual at most $a$ must have oscillation at least $(t-a)_+/t^k$. Conversely choose an oscillation $(t-a)_+/t^k$ between the incorrect and correct labels. This attains that residual bound. Hence
\[
 \calC(a)=\sup_{0<t\le1}(t-a)_+/t^k,
\]
which is comparable to $a^{1-k}$ by the choices $t=2a$ for the lower and $t\ge a$ for the upper. Theorem~\ref{thm:transfer} now produces the sharp rate by balancing its two terms.

A one-label stationary program has the same decision distribution before every report and cannot retain whether a revelation occurred. At $\theta=\pm1$ its worst error is at least $1/2$, attained by fair decisions. The two-label program has average zero for each parameter, and the uniform bound already proved makes its worst-model finite-horizon risk tend to zero. No internal horizon or uncharged cycle number is used.
\end{proof}
This is a task-sensitive retained-learning law, not a claim about ordinary bandit feedback with observed losses. Finite-memory learning itself is classical~\cite{HellmanCover}; the role of the example here is to attain both the uniform corrector modulus and the all-policy causal acquisition value through a recurrent-rank change.

\begin{example}[Duration and exit cannot be separated]\label{ex:marked}
There are three boundary labels. From label 1, exit to label 2 with probability $1/2$ after one zero-score call, or to label 3 with probability $1/2$ after three zero-score calls. Labels 2 and 3 are absorbing with unit-duration cycles and respective scores zero and one. Then
\[
 g=(1/2,0,1)^\top,\quad d=(2,1,1)^\top,\quad
 T_{1\cdot}=(0,1/2,3/2).
\]
The first coordinate of $r-Tg$ is $-3/2$, whereas that of $r-Dg$ is $-1$. From label 1 the exact finite-call risk is $1/2-3/(2N)$ for $N\ge3$. Thus the duration--exit mark matters even in a three-state experiment. This is a check on the information pattern of the bias equation, not a new counterexample to classical semi-Markov theory.
\end{example}

\begin{example}[What the task criterion excludes]\label{ex:discontinuity}
Replace the rare-revelation loss amplitude $t$ by one on the incorrect label while keeping $P_t$ above and assigning the same positive incorrect-label loss at $t=0$. From that label the average is zero for $t>0$ and one at $t=0$. Finite responses remain continuous in this version because the reward no longer depends on $t$, but gain continuity fails. Approximate corrector prices are infinite for every $a<1$, and no uniform acquisition limit is asserted. Conversely constant zero rewards make every boundary degeneration invisible. These two cases show why cost-universal recurrence and task-visible regularity are genuinely different conditions.
\end{example}

\section{Two raw realizations}\label{sec:realizations}
\subsection{Controlled queue excursions}
A physical preparation places one job in a queue. At every service call the workload $W\ge1$ moves to $W+1$ with probability $p_\theta(a)$ and to $W-1$ otherwise, ending the excursion on hitting zero. Actions and finite noisy workload reports may depend on the retained observer state and past visible reports. The preparation call is charged and does not overwrite the observer label. Scores are bounded in $[0,1]$. Assume compact parameter sets, continuous raw report probabilities and action probabilities, and
\begin{equation}\label{eq:queue-drift}
 0\le p_\theta(a)\le\bar p<1/2
\end{equation}
for every legal action. Finite report kernels can be arbitrary functions of workload subject to this finite-response continuity; they do not alter the killed workload drift.

Fix a directed graph on the $m$ observer labels. Updates are allowed only along its reflexive reachability relation. On a first-service-call end, every designated graph edge has conditional update probability at least $\zeta>0$. The required row floors must be feasible: $k_i\zeta\le1$ for every row with $k_i$ required edges. These are a compact support-stable program architecture, not all possible queue controllers.

\begin{theorem}[Unbounded physical queue realization]\label{thm:queue}
The queue architecture above satisfies the continuous-response and uniform-return hypotheses of Theorem~\ref{thm:transfer}. It has a primitive, program-uniform group-inverse bound and therefore satisfies the task criterion for every bounded continuous scored task. Robust attainment, finite-model determination and the same-memory physical-call and discount bounds hold on this entire architecture. A certified approximation can truncate the physical workload using the explicit exponential tail below. No finite-dimensional filter or geometry of the full workload posterior is assumed.
\end{theorem}
\begin{proof}
Choose $c=\frac12\log((1-\bar p)/\bar p)$ when $\bar p>0$; the case $\bar p=0$ has deterministic immediate service termination. With $V(w)=e^{cw}$, the killed transition satisfies, uniformly over history-adaptive actions,
\[
 \E[V(W_{n+1})\one_{\{W_{n+1}>0\}}\mid W_n=w,a]
 \le\rho V(w),\qquad \rho=2\sqrt{\bar p(1-\bar p)}<1.
\]
Induction gives $\Pp(T_0>n)\le e^c\rho^n$ for the service hitting time from one. Adding the single preparation call preserves a common exponential return envelope, independent of entering label and controller. Finite-call workloads are bounded by the call count plus one, so only finitely many raw probabilities enter each response. Their continuity and Lemma~\ref{lem:response} give finite-response continuity.

The first service call ends with probability at least $1-\bar p$. Consequently each designated edge of the boundary matrix has probability at least $a_*= (1-\bar p)\zeta$. The boundary graph has exactly the same terminal strongly connected components as the designated graph: it contains every designated edge, and added edges can only follow an already existing directed path. Put $L=\max(1,m-1)$. From any state there is a path of at most $L$ edges to one selected root in a reachable terminal component. In each block of $L$ boundary transitions, the probability of hitting the set of roots is at least $a_*^L$. Thus the maximal expected boundary count to that set is at most $H_*=L/a_*^L$.

For completeness, a root-hitting bound gives $\norm{G_P}_\infty\le4H_*$. For a bounded reward $f$, center it by $b=f-\Pi_Pf$ and stop its sum at the selected root. Its expected absolute sum is at most $2H_*\norm f_\infty$. In a recurrent component the stationary centering makes the mean sum over a root-to-root return zero, so this stopped sum solves $(I-P)v=b$ also at the roots; transient states satisfy the equation by first-step conditioning. Subtracting $\Pi_Pv$ identifies $G_Pf$ and gives the factor four. Corollaries~\ref{cor:group} and~\ref{cor:robust} now apply. The bound is conservative and can be replaced by direct low-dimensional evaluation when available. Truncation errors are controlled by the common exponential return tail, not by discarding rare busy periods at zero cost.
\end{proof}
This realization verifies a genuinely unbounded physical experiment. It does not prove a full predictive-dimension law for the queue workload. A separate finite geometric task can be attached and analyzed by Theorem~\ref{thm:geometry}, but that task must not be mislabeled as the entire workload quotient. Nor does this architecture theorem claim an unrestricted solution of average-cost queue control.

\subsection{A noncommuting rank-changing quantum instrument}
Let the physical Hilbert space be $\mathbb C v\oplus\mathbb C^2$, with $v$ orthogonal to the qubit sector. Choose unit qubit vectors $\psi_+,\psi_-$ with real overlap $c_0\in(0,1)$. The unknown fixed parameter is $\theta\in[-1,1]$ and the preparation is
\begin{equation}\label{eq:quantum-prep}
 \rho_\theta=(1-t^k)|v\rangle\langle v|
               +t^k|\psi_s\rangle\langle\psi_s|,
 \qquad t=|\theta|,
\end{equation}
with $\rho_0=|v\rangle\langle v|$. Each preparation is a real physical call. The next call performs the fixed three-outcome instrument with effects
\begin{equation}\label{eq:USD}
 E_+=\frac{|\psi_-^\perp\rangle\langle\psi_-^\perp|}{1+c_0},\qquad
 E_-=\frac{|\psi_+^\perp\rangle\langle\psi_+^\perp|}{1+c_0},\qquad
 E_\bot=I-E_+-E_-.
\end{equation}
The first two effects vanish on the vacuum sector. One may use the Kraus maps $\sqrt{E_z}\rho\sqrt{E_z}$; the next preparation is fresh. The diagnostic decision and score are as in Section~\ref{sec:sharpness}, with the decision due before the instrument report. The physical phase is a visible protocol type, not an internal counter supplied to the observer.

\begin{theorem}[Noncommuting singular realization]\label{thm:quantum}
This raw instrument realizes the diagnostic theorem with $\lambda(t)=(1-c_0)t^k$. Its physical excursions have duration two, retain the observer label and satisfy the task-visible criterion through $\theta=0$, where the preparation rank and recurrent rank change. It has no uniform boundary group-inverse bound. The exact all-history diagnostic value per physical call at $N=2n$ is $D_n(\lambda)/2$, and the same two-label program attains it. The physical-time rate is $N^{-1/k}$. The construction uses noncommuting physical states and a classical finite observer; it is not a theorem about arbitrary quantum memory or arbitrary additional measurements.
\end{theorem}
\begin{proof}
Write $\psi_\pm=\cos a\,e_0\pm\sin a\,e_1$, with $0<a<\pi/4$ and $c_0=\cos(2a)$. In the qubit sector $E_++E_-=\diag(\tan^2 a,1)\le I$, proving positivity and normalization. There is no false conclusive outcome, and the correct conclusive probability on $\psi_s$ is
\[
 \frac{1-c_0^2}{1+c_0}=1-c_0.
\]
The vacuum always produces $\bot$. Hence the complete visible kernel has exactly the claimed revelation probability. For $t>0$ the two state preparations corresponding to opposite signs do not commute: their qubit rank-one projections have overlap strictly between zero and one. For $0<t<1$ the preparation has rank two, at $t=0$ rank one, and at $t=1$ it is again rank one. These are actual state and support changes, not formal coordinate singularities.

All raw matrix entries and scored responses are continuous at zero. The report alphabet is finite, the physical duration is exactly two, and the two-label gain is identically zero from both entering labels. Thus Theorem~\ref{thm:transfer} applies without a recurrence sublevel. The calculation in Theorem~\ref{thm:diagnostic}, with the constant factor $1-c_0$, gives both the minimax formula and the corrector rate. Every preparation and diagnostic measurement is charged, so $N=2n$ and the score per physical call is half the per-diagnostic score. Odd horizons differ by at most one bounded call. The standard unambiguous-discrimination mechanism~\cite{ClarkeCheflesBarnettRiis} is not claimed as new; it provides an explicit noncommuting primitive realization of the task-visible theorem.
\end{proof}
With $\cos a=\sqrt3/2$ and $\sin a=1/2$, all primitive coefficients are algebraic and Theorem~\ref{thm:algebraic} also applies. In contrast, Theorem~\ref{thm:schedule}'s common-support hypothesis fails at zero. This difference is intentional: task compatibility is the mechanism that passes the singular endpoint.

\section{Common-task composition and a matched resource law}\label{sec:composition}
\subsection{General defect transfer}
We use $\TV(\mu,\nu)=\sup_A|\mu(A)-\nu(A)|$. A bounded $[0,1]$ score changes in expectation by at most this distance; a stochastic row changes in $\ell^1$ by at most twice it. Both conventions must not be used interchangeably.

\begin{proposition}[Task-sensitive causal comparison]\label{prop:defect}
Consider two experiments and a paired class of parameter-independent programs, including any simulator's private state. Assume both complete scored processes have task-compatible finite-call bridges with errors $\omega_N$ and $\omega'_N$. Suppose the simulator supplies all virtual reports before their target decisions, preserves the common action and audit interface, and compares the complete first $N$ scored physical calls with total-variation error at most $e_N$. Then
\begin{equation}\label{eq:task-defect}
 |g-g'|\le\inf_{N\ge1}\{\omega_N+e_N+\omega'_N\}.
\end{equation}
A task readout discrepancy $\zeta_N$ adds $\zeta_N$ to the braces. A per-call complete-kernel discrepancy $\eta$ gives $e_N\le\min(1,N\eta)$ by telescoping. These assertions require the stated bridge at both endpoints, on the joint predictor--simulator state. A one-way causal simulator gives only the corresponding one-way optimal-risk comparison unless a reverse simulator is also supplied.
\end{proposition}
\begin{proof}
Insert the two finite-call expected risks between the gains. Their difference is at most $e_N$ for a common $[0,1]$ score, with any declared readout discrepancy added. Each gain-to-finite-call difference is bounded at its own endpoint. Infimizing gives~\eqref{eq:task-defect}. Coupling until first discrepancy, or telescoping signed kernels for a complete marked process, gives the per-call bound. A simulator using $s$ labels and a target predictor using $m$ labels is implemented by their joint $sm$-label state, with its actual clock and physical-call schedule. Evaluating recurrence or task compatibility only on the predictor coordinate would omit a required hypothesis.
\end{proof}
This is a clock-aligned comparison. Different physical clock maps require their own duration normalization; equality of a virtual call count does not make physical-time gains equal.

For quantum primitive comparisons the telescoping is performed on unnormalized marked completely positive maps in trace norm, not by coupling noncommuting states or dividing by a zero-probability report. A supplied Kraus-coordinate approximation can imply an instrument-norm certificate, but the resulting operational bound is independent of the Kraus representation. Report marginals without the duration, action and audit interface do not establish~\eqref{eq:task-defect}.

Equation~\eqref{eq:task-defect} allows amplification weaker than a uniform group-inverse estimate. For instance if both bridges are $O(N^{-a})$ and $e_N\le N\eta$, it gives $O(\eta^{a/(a+1)})$ by balancing. This is an upper modulus, not an automatically attained lower bound. The next construction supplies actual alternatives and a matching law rather than treating an allowed error tolerance as an irreducible loss.

\subsection{A single fully charged scored protocol}
Let $V$ have a fixed preparation law $\rho$ on a compact $K\subset\R^q$, independently in every round and independently of $\theta$. Let $\sigma\in\{-1,1\}$ be a fixed inaccessible calibration sign, not redrawn along the path. Let $0\le\delta\le1$ and $0\le e\le1/2$ be known magnitudes. Each round has four physical calls with visible types:
\begin{enumerate}
\item prepare fresh physical diagnostic data and a fresh mark $V$, retaining the observer label;
\item choose the diagnostic decision, score it, and then reveal the diagnostic report with intensity $\lambda_e(t)=(t^k-e)_+$;
\item report $V$ and update the retained label;
\item produce the fixed label readout $(c,b)\in\conv K\times[-1,1]$ before the audit and score $\norm{V-c}^2+(b-\sigma\delta)^2$.
\end{enumerate}
The diagnostic score is $t\one_{\{u\ne s\}}$ as before. The environment-owned audits in calls 2 and 4 are never feedback. Their alphabets and the early diagnostic-report deadline are part of the common measurable task interface. A simulator cannot replace the ground-truth audit by a fabricated one. All preparation, load and audit calls are charged; the physical acquisition budget is $N=4n$. Define $\calR$ to be expected total score divided by $n$ rounds. The usual physical-call average is $\calR/4$; dividing further by $D_K^2+5$ gives a $[0,1]$ normalization for each round's total score. Neither rescaling changes the resource exponents.

\begin{theorem}[Matched acquisition--memory--calibration--defect law]\label{thm:joint}
For every $n\ge1$, $M\ge2$, $0\le\delta\le1$ and $0\le e\le1/2$, the minimax value over all legal retained $M$-label programs in this protocol satisfies
\begin{equation}\label{eq:joint-sandwich}
 q_M(\rho)+\delta^2+D_n(\lambda_e)
 \le\calR(4n,M,\delta,e)
 \le q_{\lfloor M/2\rfloor}(\rho)+\delta^2+D_n(\lambda_e).
\end{equation}
The upper is realized by one stationary program using $2\lfloor M/2\rfloor$ labels. The lower diagnostic term holds even against full-history strategies. Uniformly over the displayed range,
\begin{equation}\label{eq:diagnostic-erasure-rate}
 D_n(\lambda_e)\asymp_k n^{-1/k}+e^{1/k}.
\end{equation}
Let $\varepsilon_*(e)$ denote source-to-target deficiency, uniformly in $\theta$, for one diagnostic interface with source intensity $\lambda_e$, target intensity $t^k$, the unchanged environment-owned audit, and the same early-report deadline. Then
\begin{equation}\label{eq:actual-deficiency}
 e/2\le\varepsilon_*(e)\le e.
\end{equation}
If the actual mark law has matching global covering and positive-mass small-ball dimension $d$, then
\begin{equation}\label{eq:joint-law}
 \calR(N,M,\delta,e)
 \asymp M^{-2/d}+N^{-1/k}+\delta^2+\varepsilon_*(e)^{1/k},\qquad N=4n.
\end{equation}
Constants depend on the displayed raw geometric data and $k$, not on $N,M,\delta,e$. No ambient dimension is substituted for that acquired dimension.
\end{theorem}
\begin{proof}
A program has at most $M$ geometric readouts. For every fresh $V$, irrespective of its earlier retained information, its geometric loss is at least the squared distance to that codebook; averaging gives $q_M(\rho)$. Average the two fixed calibration models $\sigma=\pm1$. Since their audit is never feedback and their visible experiments agree, the expected calibration loss is $\E b^2+\delta^2\ge\delta^2$. Independently average the two diagnostic models $\theta=\pm t$. Fresh geometric reports do not reveal their sign. The indistinguishable-before-revelation argument in Theorem~\ref{thm:diagnostic} therefore gives $D_n(\lambda_e)$. These lower terms hold simultaneously under the product of the two binary minimax priors, with $t$ chosen arbitrarily close to its supremum. They may therefore be added before taking the worst model.

For the upper, store a sign label and one of $J=\lfloor M/2\rfloor$ nearest geometric centers. Initialize the sign fairly. Retain it on $\bot$, update it on conclusive reports, and update only the center coordinate on a fresh load. Set $b=0$. The known physical protocol types select the appropriate row; no internal round count is provided. This uses $2J\le M$ labels and gives exactly the upper expression, or arbitrarily closely if a codebook infimum were not attained. Here compact finite-dimensional support does give attainment. Unused labels are assigned the same legal updates as a used label; no extra inaccessible information is stored there.

For~\eqref{eq:diagnostic-erasure-rate}, choosing $t=n^{-1/k}$ gives the $n^{-1/k}$ lower because $\lambda_e(t)\le1/n$; choosing $t=e^{1/k}$ gives the $e^{1/k}/2$ lower because no report is ever conclusive. The maximum of the two is at least half their sum. For the upper, if $t^k\le2e$ the value is at most $t/2\le(2e)^{1/k}/2$. Otherwise $\lambda_e(t)\ge t^k/2$ and the geometric sum is bounded by $\min(n,2/t^k)$, giving a constant times $n^{-1/k}$. This proves the uniform comparison, including $e=0$.

The identity simulator has complete report discrepancy $t^k-(t^k-e)_+\le e$, proving the upper in~\eqref{eq:actual-deficiency}. At $t=e^{1/k}$ the source produces only $\bot$ under both signs. Before the required target report the simulator has the same information in these two models. The two target report laws are $\mu_+=(1-e)\delta_\bot+e\delta_+$ and $\mu_-=(1-e)\delta_\bot+e\delta_-$; their total-variation distance is $e$. Any common simulated law has distance at least $e/2$ from one by the triangle inequality. Additional independent mark data and unchanged audits unavailable before the deadline do not alter this argument. Finally Lemma~\ref{lem:quantization}, the global cover, and $\lfloor M/2\rfloor\ge M/3$ for $M\ge2$ compare the two quantization terms and prove~\eqref{eq:joint-law}.
\end{proof}
For $e>0$ the blind threshold can make the individual diagnostic gain discontinuous. The sharp law above is obtained by a direct minimax argument; it does \emph{not} apply Proposition~\ref{prop:defect} while silently dropping an endpoint regularity condition. The quantity $\varepsilon_*$ is the deficiency of a real pair of raw interfaces and is bounded, not falsely asserted to equal the allowed symbol $e$.

\begin{corollary}[Singular acquired marks]\label{cor:singular-marks}
Theorem~\ref{thm:joint} applies to the standard middle-third Cantor preparation, with $d=\log2/\log3$, and to products of such measures and interval measures, with their summed acquired dimensions. More generally a finite mixture with positive weight on a highest-dimensional Ahlfors-regular component and global covering of that order has the same exponent, retaining that component's actual weight in the lower constant.
\end{corollary}
\begin{proof}
A depth-$j$ Cantor cylinder has mass $2^{-j}$ and diameter $3^{-j}$. Choosing $j$ so that $3^{-j}$ is comparable to $r$ shows that an interval of radius $r$ intersects at most a fixed number of depth-$j$ cylinders, and hence has mass at most $Cr^{\log2/\log3}$. The $2^j$ cylinders supply a global cover at that scale. Products preserve these estimates by enclosing Euclidean balls in product intervals and using product covers. For a finite mixture, keep the highest-dimensional component as the submeasure in~\eqref{eq:smallball}; cover the union, using that lower-dimensional components require no larger order. These are actual preparation and scored occupation measures of the four-call protocol, not conditional laws obtained by discarding low-probability outcomes.
\end{proof}

\subsection{Separate computational resources}
If geometric centers and readouts are represented to error $\xi$, the corresponding upper root distortion is at most $\sqrt{q_J(\rho)}+\xi$, so its squared term is at most $(\sqrt{q_J(\rho)}+\xi)^2$. In the general recursion the numerical term appears in~\eqref{eq:geometric-upper}. A particular finite implementation must supply terminating arithmetic or integration certificates for its row probabilities and decisions. A $p$-bit description of $J$ centers in $\R^q$ uses $O(Jqp)$ read-only bits, a nearest-center scan can use $O(qp+\log J)$ temporary bits, and it performs $O(Jq)$ fixed-precision distance operations per load. The retained sign--center coordinate uses $\lceil\log_2(2J)\rceil$ bits. These are explicit upper counts for the stated scan, not simultaneous optimal lower bounds on all resources. Continuous raw report reading and random-bit sampling remain part of the declared acquisition interface; their finite-precision implementation must be charged rather than assumed free.

\section{Scope of the general conclusion}\label{sec:scope}
The main theorem is an equivalence for a compact continuous-response family with uniformly integrable physical returns. It is neither a declaration that every experiment has such returns nor a requirement that every slow boundary mode be uniformly controlled for every reward. Its exact obstruction is task compatibility at limiting ergodic projections. The algebraic theorem makes that obstruction decidable on finite algebraic instruments; the separate recurrence schedule gives an operational near-optimal family when the average objective is not continuous. These are distinct conclusions with different hypotheses.

The geometric conclusion concerns actual scored occupation. A finite acquired dimension gives a causal rate only after global candidate-valid covering, actual-law stability, a common executable initialization when needed, and state accounting have been established. Optimizing the controller requires the uniform lower witness and a fixed-size attaining exploration in Theorem~\ref{thm:geometry}. The four-call experiment verifies these optimized hypotheses directly for its geometric factor and derives its diagnostic lower against all history-dependent strategies. The queue realization does not claim the same dimension theorem for the entire workload posterior.

Physical regeneration is a substantive raw assumption. It does not mean observer overwrite, and it is not inferred from an arbitrary return decomposition. Common-program compactness here uses finite-product response domination or an explicit direct construction; general nondominated adaptive families remain outside that proof. The prediction quotient also has separate standard-Borel realization assumptions. Positivity of kernels alone supplies none of these extra properties.

Several natural problems remain. A necessary-and-sufficient recursive geometric classification beyond the cover-and-stability certificate is not proved. The efficient synthesis complexity of the global algebraic certification is not classified. Nor are the optimal joint costs of description length, workspace, physical computation time, numerical precision and simulator memory. Arbitrary unbounded stopping rules, nonregenerative infinite physical histories and unbounded quantum operators require further arguments. No total-variation estimate in this article is promoted to a large-deviation equivalence.

The present results isolate a common mechanism nevertheless: only the dynamically persistent directions visible to executable tasks need be paid for in a task-specific causal transfer. The marked approximate-corrector price quantifies that payment, acquired occupation identifies the spatial distortion, and an executable program must pay separately for preserving and updating both. This connects the original experimental foundation to a matching finite-resource law without replacing it by one realization.

\begin{thebibliography}{99}
\bibitem{BasuPollackRoy}
S. Basu, R. Pollack and M.-F. Roy, \emph{Algorithms in Real Algebraic Geometry}, second edition, Algorithms and Computation in Mathematics 10, Springer, Berlin, 2006. doi:10.1007/3-540-33099-2.
\bibitem{BewleyKohlberg}
T. Bewley and E. Kohlberg, The asymptotic theory of stochastic games, \emph{Math. Oper. Res.} \textbf{1} (1976), 197--208. doi:10.1287/moor.1.3.197.
\bibitem{BewleyKohlbergStationary}
T. Bewley and E. Kohlberg, On stochastic games with stationary optimal strategies, \emph{Math. Oper. Res.} \textbf{3} (1978), 104--125. doi:10.1287/moor.3.2.104.
\bibitem{Blackwell}
D. Blackwell, Equivalent comparisons of experiments, \emph{Ann. Math. Statist.} \textbf{24} (1953), 265--272. doi:10.1214/aoms/1177729032.
\bibitem{ClarkeCheflesBarnettRiis}
R. B. M. Clarke, A. Chefles, S. M. Barnett and E. Riis, Experimental demonstration of optimal unambiguous state discrimination, \emph{Phys. Rev. A} \textbf{63} (2001), 040305. doi:10.1103/PhysRevA.63.040305.
\bibitem{DenardoFox}
E. V. Denardo and B. L. Fox, Multichain Markov renewal programs, \emph{SIAM J. Appl. Math.} \textbf{16} (1968), 468--487. doi:10.1137/0116038.
\bibitem{FedergrunSchweitzer}
A. Federgrun and P. J. Schweitzer, A fixed point approach to undiscounted Markov renewal programs, \emph{SIAM J. Algebraic Discrete Methods} \textbf{5} (1984), 539--550. doi:10.1137/0605052.
\bibitem{Fritz}
T. Fritz, A synthetic approach to Markov kernels, conditional independence and theorems on sufficient statistics, \emph{Adv. Math.} \textbf{370} (2020), 107239. doi:10.1016/j.aim.2020.107239.
\bibitem{GrafLuschgy}
S. Graf and H. Luschgy, \emph{Foundations of Quantization for Probability Distributions}, Lecture Notes in Mathematics 1730, Springer, Berlin, 2000. doi:10.1007/BFb0103945.
\bibitem{HellmanCover}
M. E. Hellman and T. M. Cover, Learning with finite memory, \emph{Ann. Math. Statist.} \textbf{41} (1970), 765--782. doi:10.1214/aoms/1177696958.
\bibitem{LamondPuterman}
B. F. Lamond and M. L. Puterman, Generalized inverses in discrete time Markov decision processes, \emph{SIAM J. Matrix Anal. Appl.} \textbf{10} (1989), 118--134. doi:10.1137/0610009.
\bibitem{Meyer}
C. D. Meyer, Jr., The role of the group generalized inverse in the theory of finite Markov chains, \emph{SIAM Rev.} \textbf{17} (1975), 443--464. doi:10.1137/1017044.
\bibitem{ShaliziCrutchfield}
C. R. Shalizi and J. P. Crutchfield, Computational mechanics: Pattern and prediction, structure and simplicity, \emph{J. Stat. Phys.} \textbf{104} (2001), 817--879. doi:10.1023/A:1010388907793.
\end{thebibliography}

\end{document}
